% =====================================================================
%  VERSUCHSDURCHFÜHRUNG  (Abschnitt 1.3 der Anleitung)
%  Aufbau skizzieren, Durchführung kurz aber genau schildern,
%  ALLE vorgenommenen Einstellungen protokollieren!
% =====================================================================
\section{Versuchsdurchführung}
\label{cap:Versuchsdurchfuehrung}

\subsection{Versuchsaufbau}
\label{cap:Versuchsaufbau}
Beschreibung und Skizze/Foto des experimentellen Aufbaus. Fremde Abbildungen sind mit Quellenangabe zu kennzeichnen.

\begin{figure}[H]
    \centering
    %\includegraphics[width=1\linewidth]{Bilder/aufbau.png}
    \caption{Versuchsaufbau mit allen relevanten Komponenten}
    \label{fig:Aufbau}
\end{figure}

\subsection{Verwendete Geräte}
\label{cap:Geraete}
Verwendete Geräte und Materialien inkl. Inventarnummern und Einstellungen (z.B. Verstärkung, Zeitbase, Messbereiche):
\begin{table}[H]
    \centering
    \caption{Verwendete Geräte und wichtige Einstellungen}
    \begin{tabular}{c|c|c|c}
        \label{tab:Geraete}
        Gerät & Typ/Hersteller & Inventarnummer & Einstellung \\
        \hline
        Oszilloskop & Mustertyp A & 12345 & 2\,ms/div \\
        Multimeter & Mustertyp B & 67890 & Bereich 10\,V \\
        Spannungsquelle & Mustertyp C & 11121 & 5.0\,V \\
    \end{tabular}
\end{table}

\subsection{Durchführung}
\label{cap:Durchfuehrung}
Kurze, aber genaue Schilderung des Ablaufs:
\begin{itemize}
    \item Vorbereitung und Justage (z.B. Nullabgleich, Kalibrierung)
    \item Messung: Was wurde wie und wie oft gemessen ($N$ Wiederholungen)?
    \item Alle vorgenommenen Einstellungen und Änderungen notieren
\end{itemize}
Hinweis zur Schreibweise: neutrale Erzählperspektive (keine Ich-Form), kurz und prägnant.
